\documentclass[10pt, aspectratio=169]{beamer}

% --- THÈME ET POLICES ---
\usetheme{Boadilla}
\usepackage[utf8]{inputenc}
\usepackage[T1]{fontenc}
\usepackage[french]{babel}
\usepackage{graphicx}

% --- COULEURS PERSONNALISÉES ---
\definecolor{mainText}{HTML}{2d3748}
\definecolor{mainBlue}{HTML}{AED6F1}
\definecolor{strongBlue}{HTML}{2b6cb0}
\definecolor{pastelRose}{HTML}{F4C7D8}
\definecolor{pastelJaune}{HTML}{F9E79F}

% --- APPLICATION DES COULEURS ---
\setbeamercolor{title}{fg=mainText}
\setbeamercolor{frametitle}{bg=mainBlue, fg=mainText}
\setbeamercolor{palette primary}{bg=mainText}
\setbeamercolor{structure}{fg=strongBlue} % Couleur des titres de section, etc.

% Redéfinition de \textbf pour utiliser la couleur "strong"
\let\oldtextbf\textbf
\renewcommand{\textbf}[1]{{\color{strongBlue}\oldtextbf{#1}}}


%--- METADONNÉES ---
\title{Conception UX/UI - Application Cyber-Harcèlement}
\subtitle{Justification des choix graphiques et ergonomiques}
\author{Cible : Enfants 6-11 ans}
\institute{Briffard Anna, Dupouët Camille, Gadille Arnaud, Mignard Maël, Trichard Joan}
\date{\today}

\begin{document}

% --- SLIDE DE TITRE ---
\begin{frame}
  \titlepage
\end{frame}

% --- INTRODUCTION ---
\section{Introduction du Projet}
\begin{frame}{Contexte et Objectifs}
  \begin{block}{Contexte du projet}
    Développement d'une application de sensibilisation au cyber-harcèlement dans le cadre du BTS SIO.
  \end{block}
  
  \begin{block}{Objectif Principal}
    Créer une interface \textbf{rassurante et ludique} pour libérer la parole des enfants (6-11 ans) et éduquer sans traumatiser.
  \end{block}

  \begin{alertblock}{Problématique UX}
    Comment traiter un sujet grave (harcèlement) avec une cible jeune sans générer d'anxiété ?
  \end{alertblock}
\end{frame}

% --- CHARTE GRAPHIQUE ---
\section{La Charte Graphique}
\begin{frame}{Psychologie des Couleurs}
  \begin{block}{Une palette Pastel}
    Choix d'une palette douce pour réduire la charge cognitive et l'agressivité visuelle.
  \end{block}
  
  \begin{columns}[T]
    \begin{column}{.33\textwidth}
      \begin{beamercolorbox}[sep=1em, wd=\textwidth,
      rounded=true, shadow=true]{block body,bg=pastelRose}
        \centering
        \textbf{Confiance}\\ (Rose)
      \end{beamercolorbox}
    \end{column}
    \begin{column}{.33\textwidth}
      \begin{beamercolorbox}[sep=1em, wd=\textwidth,
      rounded=true, shadow=true]{block body,bg=pastelJaune}
        \centering
        \textbf{Éveil}\\ (Jaune)
      \end{beamercolorbox}
    \end{column}
    \begin{column}{.33\textwidth}
      \begin{beamercolorbox}[sep=1em, wd=\textwidth,
      rounded=true, shadow=true]{block body,bg=mainBlue}
        \centering
        \textbf{Sécurité}\\ (Bleu)
      \end{beamercolorbox}
    \end{column}
  \end{columns}
\end{frame}

% --- NAVIGATION ---
\section{Navigation : Le concept du "Hub"}
\begin{frame}{Le menu "Fleur"}
    \begin{block}{Interface}
        Menu principal en forme de fleur où chaque "pétale" est une option.
    \end{block}
    
    \begin{alertblock}{Justification UX}
        \begin{itemize}
            \item \textbf{Réduction du temps de décision :} Toutes les options sont visibles immédiatement.
            \item \textbf{Pas de menus cachés :} Adapté aux enfants qui ont une maturité numérique limitée.
            \item \textbf{Zones de clic larges :} Les "pétales" facilitent l'usage tactile.
        \end{itemize}
    \end{alertblock}
\end{frame}

% --- ECRANS ---
\section{Analyse des Écrans}
\begin{frame}{Écrans d'Activité et Contenu}
    \begin{columns}[T]
        \begin{column}{.5\textwidth}
            \begin{block}{Le Jeu (Ludification)}
                \begin{itemize}
                    \item \textbf{Feedback immédiat :} Valide les réponses avec des couleurs douces.
                    \item \textbf{Engagement :} Fait passer des messages de prévention de manière ludique.
                \end{itemize}
            \end{block}
        \end{column}
        \begin{column}{.5\textwidth}
            \begin{block}{La Lecture (Information)}
                \begin{itemize}
                    \item \textbf{Hiérarchie en Z :} Image -> Titre -> Texte court.
                    \item \textbf{Charge mentale réduite :} Un seul message clé par écran.
                    \item \textbf{Typographie lisible :} Police sans-serif adaptée aux jeunes lecteurs.
                \end{itemize}
            \end{block}
        \end{column}
    \end{columns}
\end{frame}

\begin{frame}{Écran d'Aide (Signalement)}
    \begin{alertblock}{Analyse UX (Design Émotionnel)}
        \begin{itemize}
            \item \textbf{Rassurance :} Interface très épurée pour ne pas distraire l'enfant.
            \item \textbf{Wording (Mots choisis) :} Vocabulaire simple et direct ("Je veux parler") au lieu de termes techniques ("Formulaire de contact").
            \item \textbf{Absence de Rouge :} On évite le code couleur "Danger" pour encourager la demande d'aide sans peur.
        \end{itemize}
    \end{alertblock}
\end{frame}

% --- CONCLUSION ---
\section{Conclusion}
\begin{frame}{Synthèse}
    \begin{block}{Deux impératifs}
        L'architecture de l'application et les choix de design répondent à :
        \begin{itemize}
            \item \textbf{Technique :} Une structure modulaire simple et maintenable.
            \item \textbf{Utilisateur :} Une interface bienveillante ("Care Design") qui transforme une procédure de signalement en une expérience sécurisante.
        \end{itemize}
    \end{block}
\end{frame}

\end{document}
